\chapter{Deployment Path and Verification Discipline}
\label{ch:deployment-path-verification}

The battery thesis closes with a deployment path that is operationally
serious but still evidence-bounded.  The goal is not to announce deployment.
The goal is to specify how battery deployment would proceed without breaking
the claim discipline of the rest of the book.

\section{Verification Before Expansion}

The correct next sequence is:

\begin{enumerate}
  \item preserve the locked battery core claims and manifests,
  \item regenerate every late-chapter artifact through named wrapper scripts,
  \item run manuscript validators and classify failures honestly, and
  \item only then widen the runtime path toward stronger HIL and physical
    deployment.
\end{enumerate}

\section{Battery Deployment Path}

For the battery domain, the practical deployment ladder is:

\begin{enumerate}
  \item benchmark replay with observed versus true trajectory separation,
  \item controller-in-the-loop runtime rehearsal with complete certificates,
  \item hardware-in-the-loop with real telemetry transport and relay behavior,
  \item supervised pilot operation with bounded fallback policies, and
  \item wider deployment only after those layers remain manifest-traceable.
\end{enumerate}

\section{Verification Discipline}

Every stage should preserve four invariants:

\begin{itemize}
  \item a named script path that can regenerate the output,
  \item a run-scoped artifact with stable location,
  \item a manifest entry tying the artifact to a chapter claim, and
  \item an explicit statement of what the result does \emph{not} prove.
\end{itemize}

This discipline is part of the battery contribution.  The thesis is not only
about a safer dispatch policy; it is also about a verifiable way to connect
runtime safety claims to evidence.

\section{Summary}

The battery deployment path is therefore incremental and auditable.  Each new
layer must inherit the same theorem-to-code-to-artifact traceability that
governs the current battery package.
